\documentclass[11pt,a4paper,UTF8,fontset=fandol]{ctexart}

\usepackage{geometry}
\geometry{left=2.35cm,right=2.35cm,top=2.55cm,bottom=2.55cm}
\usepackage{graphicx}
\usepackage{booktabs,longtable,array}
\usepackage{amsmath,amssymb}
\usepackage{float,xcolor}
\usepackage{hyperref,url}
\urlstyle{tt}
\graphicspath{{assets/stage2_report_20261023/}}
\hypersetup{unicode=true,colorlinks=true,linkcolor=black,citecolor=black,
  urlcolor=blue,pdftitle={Stage2 RAW Video Denoising Research Report},
  pdfauthor={Stage2 DenoiseModel Project}}
\setlength{\parindent}{2em}
\setlength{\parskip}{0.30em}
\setlength{\tabcolsep}{4.5pt}
\renewcommand{\arraystretch}{1.16}
\newcommand{\dn}{\ensuremath{\,\mathrm{DN}}}
\newcommand{\db}{\ensuremath{\,\mathrm{dB}}}
\newcommand{\psnr}{\ensuremath{\mathrm{PSNR}}}
\newcommand{\rmse}{\ensuremath{\mathrm{RMSE}}}
\newcommand{\pairhead}{\makebox[0.49\textwidth]{输入}\hfill\makebox[0.49\textwidth]{方法3（EMA）}\\[0.25em]}
\newcommand{\triplehead}{\makebox[0.32\textwidth]{输入}\hfill\makebox[0.32\textwidth]{方法3}\hfill\makebox[0.32\textwidth]{方法8专用输出头}\\[0.25em]}

\title{\textbf{视频去噪研究报告}\\[0.4em]
\large 方法探索、最终方案与多数据集泛化效果}
\author{肖纬杰\quad 邸锦承}
\date{}

\begin{document}
\maketitle
\begin{abstract}
\noindent
本报告记录 Stage2 RAW 视频去噪项目从公司数据适配、静态监督构建到跨传感器验证的研究过程。针对预训练教师在公司 RAW 上噪声与细节难以兼顾的问题，项目先建立可信的时域均值参考和实测噪声合成管线，再依次考察动态监督、冻结运动专家、教师融合、传感器适配、统一微调及推理门控。最终采用方法3：以 RViDeformer-L 为基座，混合真实跨增益样本与标定噪声合成样本，使用带有效区域掩码的重建、梯度和高通损失进行全参数微调，并固定第5000步的 EMA 权重用于推理。sc635hai 另保留方法8的专用可选输出头。

效果展示覆盖公司 MIS20S1、SC465SL、sc635hai 和 sc450ai 多个数据集与拍摄条件。方法3在 MIS20S1 645x 跨增益参考评测中达到29.57\db，在未参与训练的1000x原生参考评测中达到53.27\db；在未进行目标传感器训练的 SC465SL 和 sc635hai 静态序列上，零样本 PSNR 分别提升约10.52--13.21\db 和13.59\db。报告同时展示高增益动态及 sc450ai 的实际输出，按静态定量、动态观察和专用适配三种协议解释结果。正文依次介绍方法探索、最终方案和多数据集效果，附录集中列出采纳方案的训练超参数、损失、模型与推理配置。

\vspace{0.6em}
\noindent\textbf{关键词：} RAW 视频去噪；RViDeformer；时域均值监督；跨传感器泛化；固定 ISP
\end{abstract}

\newpage
\tableofcontents
\newpage

\section{引言}

\subsection{项目起点与研究目标}

Stage1 使用预训练的 RViDeformer-L\_CRVD 教师生成伪干净目标。在公司 MIS20S1 数据上，该教师对图卡细线、棋盘和文字边缘的处理偏软，影响了下游监督质量。Stage2 因而把研究对象转向教师本身：利用现有公司 RAW 数据进行微调，兼顾噪声抑制与高频结构，并检查模型在不同传感器、增益和运动条件下的表现。

项目没有公司数据的逐像素干净动态真值。可利用的资源是低增益静态序列、同场景不同增益拍摄、平场与黑帧标定，以及少量其他传感器的真实场景流。因此，研究首先需要解决监督从哪里来，再判断相应监督能够支撑哪些模型改动。

\subsection{研究过程与报告主线}

项目从静态去噪出发。方法0至方法4依次完成数据适配、参考构建、噪声合成、模型微调和评估，形成方法3这一可用基线。随后，静态监督中的运动排除机制暴露出动态内容被过度平滑的现象，推动了方法5至方法7的运动保真探索。跨传感器测试又引出了方法8至方法10的适配与退化建模。最后，全数据集预览中的高增益动态表现促成了方法11的置信门控实验。

这些探索均以同一基线作比较，最终默认权重仍为方法3的第5000步 EMA；方法8仅保留 sc635hai 专用输出头。下文先简述各条路线的设计、结果及取舍，再完整说明最终采纳方案，最后集中展示其在多种公司数据集上的泛化效果。方法编号沿用原实验顺序；复现文件保留项目中的原始名称。

\section{相关工作与数据基础}

\subsection{技术背景}

RAW 保留了传感器的线性响应和 Bayer 采样信息，适合在噪声统计层面建立恢复模型。CRVD/RViDeNet\cite{r1}、RViDeformer\cite{r2} 和 RVRT\cite{r3} 利用多帧信息进行恢复；HDR+\cite{r10}、Deep Burst Denoising\cite{r9} 与核预测方法\cite{r17} 则体现了对齐和帧间融合对低照度成像的作用。EMVD\cite{r15}、FastDVDNet\cite{r16} 以及 DnCNN\cite{r14}、CBDNet\cite{r20} 提供了不同的时空恢复与真实噪声建模思路；NL-means\cite{r18} 和 BM3D\cite{r19} 是非局部恢复的经典参照。本项目采用官方 CRVD 权重初始化 RViDeformer-L，在公司数据上微调。

在缺少干净真值时，Noise2Noise\cite{r4} 依赖噪声观测之间的统计关系，Noise2Void\cite{r5} 与 Noise2Self\cite{r6} 通过盲点构造自监督。静态长序列还可利用时域均值降低随机噪声，但运动区域的均值会混合不同时间的内容。项目据此将静态均值监督与动态探索分开处理，不把动态均值直接当作干净目标。

噪声标定以信号相关方差模型\cite{r7}为基础，并参考 RAW 数据合成的思路\cite{r8}。实际生成器采用逐增益经验系数、通道相关高斯采样、固定图样偏置和量化。动态探针使用 DIS 光流\cite{r11}；视觉对比使用固定 ISP 和 Reinhard 色调映射\cite{r12}。这些设置在各方法之间保持一致。

\subsection{公司数据集与实验分工}

全数据集预览覆盖4种传感器的45条原始 RAW 流，其中27条为场景流、18条为标定流。按传感器统计，MIS20S1 为29条（含18条标定），SC465SL 为12条，sc635hai 和 sc450ai 各2条。包含其他算法2DNR/3DNR处理的03组派生流不计入这45条。

\begin{table}[H]
\centering\small
\caption{主要数据集、内容与实验用途。}
\label{tab:dataset}
\begin{tabular}{p{0.20\textwidth}p{0.29\textwidth}p{0.40\textwidth}}
\toprule
数据集 & 内容 & 用途 \\
\midrule
MIS20S1 暗房 & 同场景7档增益，10x--1000x & 10x静态主参考；128x/313x/632x/645x真实训练；1000x真实增益留出 \\
MIS20S1 室外 & 黄昏公园、人物运动，4档增益 & 场景迁移与高增益动态展示 \\
MIS20S1 标定 & 9档平场、9档黑帧，每流50帧 & 噪声拟合与合成验证；不作为场景泛化成绩 \\
SC465SL 暗室 & 256x、2000x，各300帧 & 未参与方法3训练的跨传感器静态定量 \\
SC465SL 超市与复杂纹理 & 增益梯子及 difficultpattern，共10流 & 多场景视觉与复杂纹理动态测试 \\
sc635hai & 131x低色温150帧、16x运动彩色200帧 & 方法3零样本测试；131x另用于方法8适配，16x保留 \\
sc450ai & 804x灯光变化150帧、1435x黑物运动210帧 & 跨传感器压力测试；1435x不参与适配 \\
\bottomrule
\end{tabular}
\end{table}

\subsection{RAW 适配与参考资格}

不同数据源的位对齐、CFA和黑电平并不统一。04场景流和05组通常需要从高12位右移4位；02组采用低位对齐。sc450ai 虽然附带说明写为高12位和4128\dn 黑电平，实测却是低位对齐、有效10位，样本均落在4\dn 格点，黑电平约256--258\dn。若直接套用说明，会把有效信号裁掉。场景流以 RGGB 为主，标定流为 GRBG；所有数据经适配器转换到统一的 R/Gr/Gb/B 通道顺序。

黑电平逐流、逐通道测量；帧数从文件字节数核实。标称增益不作为网络输入条件，因为不同档位伴随照明变化，实际噪声不随标签单调增加。例如暗房128x的分半绿色通道噪声约79.05\dn，1000x约28.93\dn。噪声合成因此使用标定实测系数，而非只按场景文件名推断强度。

监督资格随数据而定。10x低增益静态参考用于主方法训练；同流奇偶分半参考用于可信静态区域的测试或适配。旧教师输出、运动时域均值和盲点观测只作为对应探索的信号，不被登记为干净动态真值。

\subsection{评测口径}

静态评测在有效区域计算 native DN 域的 PSNR、RMSE、高频保持率和残差。跨传感器静态测试使用互斥奇偶帧池：偶数帧片段对奇数帧参考，再交换方向；只保留在全部、奇数、偶数帧池中均达到90\%静态置信度的像素。运动筛选采用 $\mathrm{median}+3\times1.4826\times\mathrm{MAD}$ 的鲁棒阈值。

动态序列没有干净逐帧目标，报告其时序响应、高频与裁剪等诊断量，并结合固定 ISP 视频观察。高频比下降既可能来自去噪，也可能来自结构损失；temporal MAD 下降既可能表示抑制噪声，也可能表示运动被平均。因此，两者均与视觉结果一起解释。

所有同序列输入与输出共享白平衡和曝光，关闭 ISP 的2DNR/3DNR。全数据集展示采用 edge-aware 去马赛克、高光重建和 Reinhard 映射。PSNR 在 RAW 域计算，展示视频用于观察降噪、纹理与运动；不同数据集的参考域和有效区域不同，绝对 PSNR 不作跨数据集排名。

\section{试过的方法与方案选择}
\label{sec:exploration}

\subsection{方法0至方法4：形成静态基线}

方法0建立逐流适配、配准、曝光映射与运动掩码。方法1用暗房10x的200帧构建静态参考，奇偶均值残差为0.51--0.80\dn，曝光归一化后的跨增益偏差为1.17--1.86\dn。方法2利用平场和黑帧生成9档合成噪声，最大逐通道统计误差8.66\%，低于15\%验收阈值。上述准备共同支撑方法3的混合数据微调，方法4负责评估。

评估中出现过一次口径修正：1000x与10x跨捕获曝光关系在高亮处偏离线性，同一模型按该映射仅有18.84\db；配准复核后，改用1000x自身的静态参考得到53.27\db，奇偶互斥验证为53.10/53.07\db。该结果使项目保留方法3，并将跨捕获映射与模型误差分开检查。方法3的监督、损失和推理流程在第\ref{sec:adopted}节统一说明。

\subsection{方法5至方法7：运动保真探索}

方法3的静态去噪改善后，645x视频中的运动人物仍有半透明残影。原运动审计中，输出的运动区帧差仅为输入的27\%--32\%。这推动了三条后续路线，比较时均要求运动响应改善且静态、原生增益及高频指标不回退。

\begin{table}[H]
\centering\small
\caption{运动保真路线的简要比较。}
\label{tab:motion_trials}
\begin{tabular}{p{0.10\textwidth}p{0.30\textwidth}p{0.38\textwidth}p{0.10\textwidth}}
\toprule
方法 & 设计 & 主要结果 & 取舍 \\
\midrule
方法5 & 光流对齐的动态 N2N；教师静态锚；冻结范围与损失权重实验 & 收紧运动置信规则后覆盖仅1.81\%--3.82\%；多组候选未同时通过运动、静态和1000x门槛 & 未采纳 \\
方法6 & 冻结方法3，训练空间残差运动专家；共4版 & 门控 IoU 达0.533，但真实动态内容未稳定恢复；更强残差伴随静态偏移或目标误差增大 & 未采纳 \\
方法7 & 在运动区融合旧教师残差，扫描36组配置 & 最强配置的运动比相对增加0.9162\%；固定 ISP 未见稳定的人物轮廓改善 & 未采纳 \\
\bottomrule
\end{tabular}
\end{table}

方法5用 DIS 光流将后帧对齐到当前帧，以前后向误差和光度误差筛选可信像素。增加严格的1像素运动门槛后，监督覆盖明显缩小。方法6转而冻结主干，减少全模型漂移；方法7再取消训练，检验旧教师残差是否能直接补充运动信息。三者依次减少了参数更新带来的影响，但在既定数据与验收条件下，均未得到足以替换基线的结果。

因此，后续实验继续保留方法3，研究重心转向跨传感器：先测同一权重在哪些数据域有效，再决定是否需要目标域适配。

\subsection{方法8至方法10：跨传感器适配与退化验证}

零样本测试表明，方法3在 SC465SL 暗室和 sc635hai 低色温静态域有明显降噪收益，而 sc450ai 黑色运动物体仍明显发软。方法8据此考察逐传感器适配：sc635hai 使用奇偶静态监督，sc450ai 使用单帧盲点监督，均混入 MIS20S1 回放。结果仅保留 sc635hai 的专用输出头，详细配置见第\ref{sec:head}节；sc450ai 的5个检查点均没有恢复高频，1435x端点裁剪还从3.28\%升至4.64\%。

方法9尝试统一微调4种传感器，只更新15,508个参数、训练100步。SC465SL和sc635hai静态指标略升，但 sc450ai 的高频比由0.126621降至0.126589，未改善黑衣与运动轮廓，故没有作为统一模型采纳。

方法10随后检查 sc450ai 的监督来源。现有数据只有两条动态场景流，没有独立暗场、平场或干净 RAW 配对。其补充退化实验比较了原独立噪声、4\dn量化独立噪声和频谱整形噪声；频谱整形虽改善804x的聚合相关误差，却未通过逐通道噪声与1435x保留集验证。该生成器未进入最终方法。

这组实验使最终方案保持清晰：默认使用同一方法3权重评测泛化，仅给 sc635hai 保留一个专用可选头。目标域的零样本结果与专用适配结果分别报告。

\subsection{方法11：高增益推理门控}

全数据集预览显示，MIS20S1室外1500x等序列仍有运动朦胧。四路诊断比较原始输入、正常时域推理、加权空间拼接和单帧重复8次的推理：改变拼接权重对主体影响很小，单帧重复则恢复了更多当前帧结构，同时带回明显彩噪。这些观察支持将时域融合视为该现象的重要来源。

方法11尝试用输入残差、梯度和光流置信度，在正常时域输出与单帧重复输出之间融合。两条代表序列的非零门控覆盖仅0.487\%和0.713\%，全局输出接近原时域分支，逐帧观察也未见稳定的主体重影改善，因此未纳入最终推理流程。

\subsection{最终方案的确定}

上述探索分别检验了训练监督、可训练模块、教师残差、传感器适配和推理融合。最终保留的方法3具有完整的监督与评测记录，静态恢复结果稳定；后续路线未在各自目标上满足替换条件。方法8的 sc635hai 输出头则以专用可选权重单独保留。接下来围绕这两项说明最终实现，并将其效果统一放到多数据集测试中展示。

\section{最终采纳的方法}
\label{sec:adopted}

\subsection{整体流程}

默认方法由四个环节组成：RAW适配、可信静态监督、真实与合成混合训练、EMA多帧推理。方法0至方法2提供的数据处理和监督构造均属于方法3的完整训练管线。

\begin{figure}[H]
\centering\small
\fbox{\parbox[c][1.5cm][c]{0.25\textwidth}{\centering 原始 RAW\\位对齐、CFA、黑电平\\配准与曝光映射}}
$\longrightarrow$
\fbox{\parbox[c][1.5cm][c]{0.25\textwidth}{\centering 静态参考与有效掩码\\实测噪声合成\\真实:合成约1:1}}
$\longrightarrow$
\fbox{\parbox[c][1.5cm][c]{0.25\textwidth}{\centering RViDeformer-L\\全参数微调5000步\\EMA权重推理}}
\caption{最终方法3的完整流程。训练数据进入统一参考域；输出在对应评测协议下恢复到 native DN 域，再使用固定 ISP 展示。}
\label{fig:adopted_pipeline}
\end{figure}

模型输入为8帧、4通道的 packed Bayer，输出对应的8帧。训练采用 $64\times64$ packed patch，即原 Bayer 的 $128\times128$ 区域；推理通过重叠分块处理完整画面。方法3不增加增益标签输入或传感器条件分支。

\subsection{真实样本与静态监督构建}

主参考来自 MIS20S1 暗房10x序列的200帧平均。参考构建包含固定图样修正，并记录运动及无效区域。参考形状为 $4\times540\times960$。奇偶分半残差在 R/Gr/Gb/B 上分别为0.511/0.797/0.763/0.560\dn；图卡高频检查确认平均后仍保留可用结构。

真实输入选用128x、313x、632x和645x。对每个输入增益，先读取其位对齐、黑电平和运动记录，再使用整数配准及逐通道仿射曝光映射进入主参考坐标和亮度域。设映射后的像素为 $x_c$，参考黑电平为 $b_c$，归一化为
\begin{equation}
 \widetilde{x}_c=\operatorname{clip}\left(\frac{x_c-b_c}{4095-b_c},0,1\right).
 \label{eq:normalization}
\end{equation}
目标为相同位置的静态参考，在8个时间位置上重复。

有效掩码同时满足参考有效、输入8帧均为静态、全部输入通道均未饱和。饱和阈值为映射域3800\dn，运动掩码阈值为0.5。只有这些位置参与损失，避免对错位运动和饱和内容施加静态均值目标。1000x真实序列不进入训练。

\subsection{实测噪声合成与混合采样}

合成分支仍以同一静态参考为干净目标，通过9档平场与黑帧建立噪声模型。信号相关方差可写为
\begin{equation}
 \sigma_c^2(s,G)=a_c\,sG+b_c^{\mathrm{read}}G^2.
\end{equation}
实际生成器保留了每档增益的经验 shot 系数 $\alpha_{c,g}$ 和 read variance $v_{c,g}$，用于处理标定中的非单调变化。令 $d$ 为固定图样偏置，$z\sim\mathcal N(0,I)$，$LL^\top=C$ 为聚合通道相关矩阵的 Cholesky 分解，则合成为
\begin{align}
 q &= y+d,\qquad s_c=\max(q_c-b_c,0),\\
 x^{\mathrm{syn}} &= \operatorname{clip}\left(
 \operatorname{round}\left[q+\sqrt{\alpha_g\odot s+v_g}\odot(Lz)\right],0,4095\right).
 \label{eq:synthetic}
\end{align}
该实现使用相关高斯采样匹配信号相关方差。合成后按1\dn格点量化，保留原始黑电平与合法范围。

合成增益标签为100、200、400、800、1600、3200、6400、12800和25600，每档生成8个独立噪声时序变体。8帧具有相同干净内容和不同噪声，真实与合成分支按约1:1混合。真实性检验比较合成与真实同档数据的高通差分噪声，最大逐通道相对误差8.66\%，通过15\%阈值。逐档参数列于附录。

\subsection{网络与优化目标}

网络采用 RViDeformer-L，严格加载官方 CRVD 权重后全参数微调。构造配置包含11组深度参数，前7组深度为4、嵌入维度84，后4组深度为2、嵌入维度108；注意力头数均为6，时空窗口为 $6\times8\times8$。实际模型共有7,505,716个可训练参数。完整构造参数见表\ref{tab:model_config}。

训练先将预测裁剪到 $[0,1]$，再计算带掩码的损失。记
\begin{equation}
 \langle z\rangle_M=\frac{\sum M\odot z}{\max(\sum M,1)},
\end{equation}
其中掩码沿通道维展开。总损失为
\begin{equation}
 \mathcal L=\mathcal L_{\mathrm{char}}+0.1\mathcal L_{\mathrm{grad}}+0.1\mathcal L_{\mathrm{hp}}.
 \label{eq:loss}
\end{equation}
Charbonnier 项\cite{r13}约束重建误差：
\begin{equation}
 \mathcal L_{\mathrm{char}}=\left\langle\sqrt{(\hat y-y)^2+10^{-6}}\right\rangle_M.
\end{equation}
若写为 $\epsilon^2$，则 $\epsilon=10^{-3}$。梯度项约束水平、垂直梯度的幅值差：
\begin{equation}
 \mathcal L_{\mathrm{grad}}=\tfrac12\left(
 \left\langle\left||D_x\hat y|-|D_xy|\right|\right\rangle_{M_x}+
 \left\langle\left||D_y\hat y|-|D_yy|\right|\right\rangle_{M_y}\right),
\end{equation}
其中 $M_x,M_y$ 是相邻有效掩码的乘积。高通项为
\begin{equation}
 H(u)=u-\operatorname{AvgPool}_{3\times3}(u),\qquad
 \mathcal L_{\mathrm{hp}}=\langle|H(\hat y)-H(y)|\rangle_M.
\end{equation}
高通仅在空间计算，不跨时间平均。重建项负责整体恢复，两个0.1权重的结构项补充边缘和局部高频约束。

\subsection{训练、权重选定与完整画面推理}

方法3使用 batch size 1、固定学习率 $10^{-4}$ 的 Adam，以FP32完成总计5000步训练。训练从第200步恢复至第5000步；总步数不额外累加。EMA衰减为0.999，逐优化步更新。最终固定第5000步 EMA，后续探索和多数据集测试均使用该权重作为默认模型。

完整画面推理采用8帧时间窗口、时间步长4，空间 tile 为 $64\times64$ packed、步长56。空间和时间重叠处等权平均；不足8帧时重复末帧补齐。输出归一化裁剪后，按相应评测管线恢复 DN 值。分块批大小仅影响一次并行处理多少 tile，不改变训练 batch size。

\subsection{sc635hai 专用可选输出头}
\label{sec:head}

方法8在方法3 EMA基础上冻结主干，仅更新 \texttt{norm}、\texttt{conv\_after\_body} 和 \texttt{conv\_last}，共6个张量、12,400个参数。对应头部由108维 LayerNorm、$108\rightarrow84$线性映射及 $84\rightarrow4$、核为 $1\times3\times3$ 的卷积组成。

适配使用 sc635hai 131x低色温序列，奇偶输入与参考帧池互斥；以0.5概率混入 MIS20S1 回放，减少原域漂移。batch size为1、学习率 $10^{-5}$、训练500步，采用BF16 autocast、FP32损失和0.999 EMA；梯度范数裁剪为1.0。实际500步中，sc635hai采样238步，回放262步。16x运动彩色序列完整保留，不参与反向传播。

该输出头通过既定自动检查并被登记为 sc635hai 专用可选权重；默认方法3继续用于统一比较。下节分别展示默认权重的零样本结果和专用头的适配结果。

\section{多数据集泛化效果展示}
\label{sec:results}

\subsection{展示范围与分组}

本节展示的是\textbf{公司多数据集、跨传感器及跨拍摄条件的泛化表现}。默认方法3只在 MIS20S1暗房真实样本及其合成噪声上训练；SC465SL、sc635hai和sc450ai均以同一权重先做零样本测试。MIS20S1室外用于检查场景变化，1000x用于检查未参与训练的真实增益。另列的方法8结果只表示 sc635hai 专用适配。

45条原始流均完成输入与方法3的固定 ISP 对比。每流取首、中、尾3个连续8帧片段，共24帧，总计1080帧。27条场景流用于观察画质，18条标定流用于噪声与数值诊断；45/45表示推理覆盖。以下按具体数据集给出静态指标和代表画面。

\subsection{MIS20S1：训练覆盖增益与未见真实增益}

\begin{table}[H]
\centering\small
\caption{MIS20S1静态结果。645x在训练覆盖内，使用10x跨增益参考；1000x真实序列未参与训练，使用原生静态参考。}
\label{tab:mis20s1}
\begin{tabular}{lrrrrrr}
\toprule
协议 & 输入PSNR & 方法3 PSNR & 提升 & 输入RMSE & 方法3 RMSE & HF比 \\
 & (dB) & (dB) & (dB) & (DN) & (DN) & \\
\midrule
645x跨增益参考 & 24.3519 & 29.5729 & +5.2210 & 248.1197 & 136.0228 & 0.9712 \\
1000x原生参考 & 44.6204 & 53.2675 & +8.6471 & 24.0565 & 8.8895 & 0.8562 \\
\bottomrule
\end{tabular}
\end{table}

645x的噪声显著下降，高频比0.9712，说明静态图卡结构与参考较为接近。图\ref{fig:mis645}展示同帧输入与最终输出：墙面和图卡周围的彩噪减少，棋盘、文字及卡片轮廓仍可分辨。由于645x参与训练，这一行用于说明训练域的恢复效果，不作为未见增益成绩。

\begin{figure}[H]
\centering\pairhead
\includegraphics[width=\textwidth,trim=0 0 0 44,clip]{fig04_a3_645x_isp.jpg}
\caption{MIS20S1暗房645x固定 ISP 对比，上为全景，下为固定中心裁剪。该图展示最终方法3在训练覆盖域的恢复效果；静态定量只统计有效区域。}
\label{fig:mis645}
\end{figure}

1000x原生参考评测的PSNR提升8.65\db，RMSE由24.06\dn降至8.89\dn。为检查同流参考的影响，另使用互斥奇偶帧池验证，方法3得到53.10/53.07\db。该结果与完整原生参考评测接近，支持模型对未参与训练的1000x真实流具有静态去噪能力；它仍是同一暗房场景下的真实增益测试。

\begin{figure}[H]
\centering
\makebox[0.32\textwidth]{输入}\hfill\makebox[0.32\textwidth]{方法3（EMA）}\hfill\makebox[0.32\textwidth]{1000x原生参考}\\[0.25em]
\includegraphics[width=\textwidth,trim=0 133 0 135,clip]{fig20_native_1000x.jpg}
\caption{MIS20S1 1000x输入、方法3与原生参考。该序列不参与主方法训练；展示保持相同亮度映射，定量按表\ref{tab:mis20s1}所列协议计算。}
\label{fig:mis1000}
\end{figure}

\subsection{SC465SL与sc635hai：跨传感器零样本静态泛化}

表\ref{tab:zero_shot}中所有结果使用同一方法3 EMA，没有目标传感器训练。SC465SL两档增益的提升为10.52--13.21\db，sc635hai低色温静态提升约13.59\db。两个奇偶方向的结果接近，且高频保持率为0.911--0.958；输出 legal RAW fraction 均为100\%。

\begin{table}[H]
\centering\small
\caption{不同数据集上的零样本静态泛化。偶/奇表示输入帧池，参考使用相反奇偶池；HF为输出/参考高频 RMS，MAD为时序平均绝对差。}
\label{tab:zero_shot}
\begin{tabular}{llrrrrr}
\toprule
数据集 & 输入池 & 输入PSNR & 方法3 PSNR & 提升 & HF比 & MAD 输入$\to$输出 \\
 & & (dB) & (dB) & (dB) & & (DN) \\
\midrule
SC465SL 256x & 偶 & 41.840 & 55.052 & +13.211 & 0.957 & 31.889$\to$1.567 \\
 & 奇 & 41.840 & 54.999 & +13.158 & 0.958 & 31.886$\to$1.578 \\
SC465SL 2000x & 偶 & 46.039 & 56.563 & +10.523 & 0.911 & 20.238$\to$1.312 \\
 & 奇 & 46.034 & 56.647 & +10.613 & 0.912 & 20.248$\to$1.301 \\
sc635hai 131x & 偶 & 38.737 & 52.325 & +13.588 & 0.919 & 47.481$\to$2.165 \\
 & 奇 & 38.736 & 52.324 & +13.587 & 0.921 & 47.481$\to$2.163 \\
\bottomrule
\end{tabular}
\end{table}

\begin{figure}[H]
\centering
\includegraphics[width=\textwidth]{generalization_static.pdf}
\caption{方法3在三组未参与训练的目标域静态数据上的泛化：左为相对输入的 PSNR 提升，右为高频保持率。数据来自表\ref{tab:zero_shot}，参考均采用互斥奇偶帧池。}
\label{fig:zero_shot_plot}
\end{figure}

SC465SL两档结果在图\ref{fig:sc465_256}和图\ref{fig:sc465_2000}中展示。输入中的细碎彩噪和背景波动得到抑制，风扇、人偶、线缆及背景卡片仍可辨。两图来自完整画面预览，静态定量则按各自的有效区域计算。

\begin{figure}[H]
\centering\pairhead
\includegraphics[width=\textwidth,trim=0 0 0 47,clip]{generalization_sc465sl_256x.png}
\caption{SC465SL暗室256x的零样本结果，上为全景，下为中心裁剪。模型未在SC465SL训练；固定 ISP 下可观察到明显降噪及背景结构保留。}
\label{fig:sc465_256}
\end{figure}

\begin{figure}[H]
\centering\pairhead
\includegraphics[width=\textwidth,trim=0 0 0 47,clip]{generalization_sc465sl_2000x.png}
\caption{SC465SL暗室2000x的零样本结果，使用与256x相同的方法3权重和推理流程。静态PSNR提升约10.52--10.61\db，高频比约0.911--0.912。}
\label{fig:sc465_2000}
\end{figure}

sc635hai 131x在低色温和较大蓝通道白平衡补偿下，仍显示较强的静态去噪能力。图\ref{fig:sc635_static}给出输入、方法3与相反奇偶池静态参考，展示背景卡片和人偶结构的保留。表\ref{tab:zero_shot}对应方法3的零样本结果；后续使用131x训练专用头所得的数字单独列在第\ref{sec:head_results}节，保持两个实验的含义独立。

\begin{figure}[H]
\centering
\makebox[0.32\textwidth]{输入}\hfill\makebox[0.32\textwidth]{方法3（零样本）}\hfill\makebox[0.32\textwidth]{相反奇偶池参考}\\[0.25em]
\includegraphics[width=\textwidth,trim=0 710 0 40,clip]{generalization_sc635hai_131x.png}
\caption{sc635hai 131x低色温数据集的零样本静态对比，选自偶输入方向的第8帧。参考为对侧奇数帧池的静态均值，动态像素不计入静态指标。}
\label{fig:sc635_static}
\end{figure}

\subsection{sc635hai专用输出头：适配域与动态保留集}
\label{sec:head_results}

方法8的131x指标来自训练数据尾段复评，该尾段参与过训练采样，不计为独立泛化结果。表\ref{tab:head_results}显示，专用头在静态域仅带来约0.034\db的小幅变化，高频与时序指标保持接近。其16x动态保留集未参加训练；由于缺少干净动态目标，只报告诊断量。

\begin{table}[H]
\centering\small
\caption{sc635hai专用头效果。131x为训练数据尾段复评，16x为完整保留的动态序列；本表与零样本表采用不同评测样本，基线数值不混用。}
\label{tab:head_results}
\begin{tabular}{llrr}
\toprule
数据与协议 & 指标 & 方法3 & 方法8专用头 \\
\midrule
131x适配域 & 静态PSNR，偶方向(dB) & 52.3535 & 52.3875 \\
 & 静态PSNR，奇方向(dB) & 52.3552 & 52.3884 \\
 & HF比（约） & 0.9199 & 0.9181 \\
 & temporal MAD(DN) & 2.1628 & 2.1247 \\
\midrule
16x动态保留集 & HF诊断比 & 0.763605 & 0.763078 \\
 & motion/input ratio & 0.599602 & 0.596438 \\
 & temporal MAD(DN) & 14.6745 & 14.6218 \\
 & 端点裁剪 & 2.1994\% & 1.9344\% \\
 & legal RAW fraction & 100\% & 100\% \\
\bottomrule
\end{tabular}
\end{table}

\begin{figure}[H]
\centering\triplehead
\includegraphics[width=\textwidth,trim=0 0 0 44,clip]{fig11_a8_sc635hai.jpg}
\caption{sc635hai 16x动态保留集：输入、默认方法3和方法8专用输出头。两种输出在结构与纹理上接近，固定 ISP 首、中、末帧复核未见明显新增重影、断边或色偏。}
\label{fig:sc635_head}
\end{figure}

方法8的变化主要集中在输出映射，未形成可辨的动态结构提升，因此保留其“sc635hai专用可选头”的身份。跨传感器静态泛化的主要证据仍是未经目标域适配的方法3。

\subsection{室外、高增益和sc450ai数据集的实际输出}

多数据集展示也包含强噪声、运动和复杂纹理。MIS20S1室外1500x与训练暗房具有场景和噪声差异。图\ref{fig:outdoor}中方法3明显压低彩噪，但人物局部有姿态平均，树木纹理偏软。该序列没有干净动态参考，因此不提供动态PSNR。

\begin{figure}[H]
\centering\pairhead
\includegraphics[width=\textwidth,trim=0 0 0 44,clip]{fig15_a3_1500x_isp.jpg}
\caption{MIS20S1室外1500x跨场景输出。上为全景，下为中心裁剪；同一默认方法3显著降噪，同时可见人物与背景细节的平滑现象。}
\label{fig:outdoor}
\end{figure}

sc450ai 1435x黑物运动在零样本诊断中，高频比为0.127，端点裁剪约3.28\%，运动/静态时序比为0.057/0.088。图\ref{fig:sc450}保留了方法3原始结果：黑衣内部结构偏空、边界偏软。804x灯光变化的高频比为0.690，暗部仍有低频色斑。SC465SL difficultpattern视频也观察到拖影和重影。上述画面与前三组静态定量共同构成了多数据集展示的完整结果。

\begin{figure}[H]
\centering\pairhead
\includegraphics[width=\textwidth,trim=0 0 600 44,clip]{fig12_a8_sc450ai.jpg}
\caption{sc450ai 1435x黑物运动的零样本结果。此处仅显示原对比图的输入与最终方法3两列；黑色衣物内部与运动轮廓仍偏软。}
\label{fig:sc450}
\end{figure}

\begin{table}[H]
\centering\small
\caption{最终方法在不同公司数据集中的展示结果。}
\label{tab:result_scope}
\begin{tabular}{p{0.26\textwidth}p{0.23\textwidth}p{0.40\textwidth}}
\toprule
数据条件 & 使用权重与协议 & 已观察结果 \\
\midrule
MIS20S1暗房645x & 方法3；训练覆盖域、跨增益参考 & +5.22\db，HF 0.971 \\
MIS20S1暗房1000x & 方法3；真实增益留出、原生参考 & +8.65\db；互斥奇偶验证53.10/53.07\db \\
SC465SL暗室256x/2000x & 方法3；跨传感器零样本静态 & +10.52--13.21\db，HF 0.911--0.958 \\
sc635hai低色温131x & 方法3；跨传感器零样本静态 & 约+13.59\db，HF 0.919--0.921 \\
sc635hai动态16x & 方法3及方法8专用头；动态保留集 & 两者画面接近；专用头未见明显新增伪影 \\
MIS20S1室外1500x & 方法3；跨场景动态观察 & 噪声下降，人物姿态与树木结构偏平滑 \\
sc450ai、复杂纹理动态 & 方法3；零样本压力测试 & 可见结构平滑、色斑或重影，不计为通过的静态泛化结果 \\
\bottomrule
\end{tabular}
\end{table}

方法3的静态恢复收益已在 MIS20S1、SC465SL和sc635hai多个数据集上得到对应协议的验证。固定 ISP 画面进一步给出了同一权重在不同场景中的实际表现；专用适配头与动态观察均按各自数据身份报告，不并入零样本静态成绩。

\clearpage
